\documentclass[11pt,a4paper]{article}
\usepackage[utf8]{inputenc}
\usepackage[T1]{fontenc}
\usepackage{lmodern}
\usepackage[french]{babel}
\usepackage{amsmath,amssymb,mathtools}
\usepackage{icomma}
\usepackage[version=4]{mhchem}
\usepackage{siunitx}
\sisetup{output-decimal-marker={,},per-mode=symbol}
\usepackage{xcolor}
\usepackage{colortbl}
\usepackage{tikz}
\usepackage[most]{tcolorbox}
\usepackage{enumitem}
\usepackage{array}
\usepackage{booktabs}
\usepackage{fancyhdr}
\usepackage[hidelinks]{hyperref}
\usetikzlibrary{arrows.meta,positioning,calc}
\usepackage[a4paper,margin=2.1cm,headheight=26pt,headsep=10pt]{geometry}

\definecolor{primary}{RGB}{146,95,161}
\definecolor{primarydark}{RGB}{92,55,108}
\definecolor{acidecol}{RGB}{205,60,45}
\definecolor{basecol}{RGB}{40,110,200}
\definecolor{methcol}{RGB}{90,90,100}

\tcbset{boxbase/.style={enhanced,breakable,boxrule=0.9pt,arc=2.5pt,left=3mm,right=3mm,top=2.3mm,bottom=2.3mm,fonttitle=\bfseries,coltitle=white,toptitle=0.6mm,bottomtitle=0.6mm}}
\newtcolorbox[auto counter]{exobox}[1][]{boxbase,colback=primary!4,colframe=primary,title={\textsc{Exercice}~\thetcbcounter\ifx\relax#1\relax\relax\else\textmd{ --- \itshape #1}\fi}}
\newtcolorbox{consigne}{boxbase,colback=methcol!7,colframe=methcol,sharp corners}

\pagestyle{fancy}
\fancyhf{}
\fancyhead[L]{\small\textcolor{primarydark}{\textbf{Fiche 10} — Thème 1 : couples acide/base}}
\fancyhead[R]{\small\textcolor{primarydark}{\textsc{Exercices — Difficile}}}
\fancyfoot[C]{\small\thepage}
\renewcommand{\headrulewidth}{0.8pt}
\renewcommand{\headrule}{\hbox to\headwidth{\color{primary}\leaders\hrule height \headrulewidth\hfill}}

\setlength{\parindent}{0pt}
\setlength{\parskip}{3pt}

\begin{document}

\begin{center}
{\LARGE\bfseries\color{primarydark} Thème 1 --- Niveau Difficile}\\[2pt]
{\color{primary}\itshape Espèces polyprotiques, ampholytes en contexte, acides aminés}
\end{center}

\begin{consigne}
\textbf{Objectif du niveau.} On enchaîne les gestes du thème sur des espèces à plusieurs protons et sur des
molécules organiques. Chaque question est décomposée en sous-questions (a, b, c) de difficulté croissante :
traitez-les dans l'ordre.
\end{consigne}

\begin{exobox}[l'acide phosphorique, un triacide]
L'acide phosphorique \ce{H3PO4} possède \textbf{trois} protons cessibles. Il intervient dans les tampons
biologiques et les boissons.
\begin{enumerate}[label=\alph*),leftmargin=1.6em,itemsep=3pt]
  \item Écrire les \textbf{trois} demi-équations successives de perte d'un \ce{H+}, de \ce{H3PO4} jusqu'à
        \ce{PO4^{3-}}. En déduire les trois couples acide/base successifs.
  \item Parmi les quatre espèces \ce{H3PO4}, \ce{H2PO4-}, \ce{HPO4^{2-}}, \ce{PO4^{3-}}, dire lesquelles sont
        des \textbf{ampholytes} et justifier.
  \item Écrire les deux réactions montrant que \ce{HPO4^{2-}} peut se comporter tantôt comme un acide, tantôt
        comme une base, vis-à-vis de l'eau.
\end{enumerate}
\end{exobox}

\begin{exobox}[la glycine, un acide aminé amphotère]
La glycine est le plus simple des acides aminés. En milieu très acide, elle existe sous la forme
entièrement protonée \ce{HOOC-CH2-NH3+}. Cette molécule porte deux sites acido-basiques : un groupe
carboxyle \ce{-COOH} (acide) et un groupe ammonium \ce{-NH3+} (acide lui aussi, dont la base conjuguée est
l'amine \ce{-NH2}).
\begin{enumerate}[label=\alph*),leftmargin=1.6em,itemsep=3pt]
  \item À partir de la forme \ce{HOOC-CH2-NH3+}, écrire la formule obtenue quand \textbf{seul} le groupe
        carboxyle perd son \ce{H+}. Cette espèce (notée \ce{^{-}OOC-CH2-NH3+}) est appelée \textbf{zwitterion}
        (elle porte à la fois une charge $+$ et une charge $-$).
  \item À partir du zwitterion, écrire la formule obtenue quand le groupe ammonium perd à son tour son
        \ce{H+}. On obtient la forme entièrement déprotonée.
  \item Montrer que le zwitterion \ce{^{-}OOC-CH2-NH3+} est un \textbf{ampholyte} : préciser lequel de ses deux
        groupes lui permet de jouer l'acide, et lequel lui permet de jouer la base.
\end{enumerate}
\end{exobox}

\begin{exobox}[lire une table et raisonner sur les couples]
On donne un extrait de table (les valeurs de $pK_a$ ne serviront qu'au repérage, pas à un calcul) :
\begin{center}\small
\begin{tabular}{lcl}
\toprule
Acide & $pK_a$ & Base conjuguée\\
\midrule
\ce{H2C2O4} (acide oxalique) & $1{,}2$ & \ce{HC2O4-}\\
\ce{HC2O4-}                  & $4{,}3$ & \ce{C2O4^{2-}}\\
\ce{CH3COOH}                 & $4{,}75$ & \ce{CH3COO-}\\
\ce{NH4+}                    & $9{,}25$ & \ce{NH3}\\
\bottomrule
\end{tabular}
\end{center}
\begin{enumerate}[label=\alph*),leftmargin=1.6em,itemsep=3pt]
  \item Dans cette table, une même espèce apparaît une fois dans la colonne « Acide » et une fois dans la
        colonne « Base conjuguée ». Laquelle ? Qu'est-ce que cela prouve à son sujet ?
  \item Donner la base conjuguée de \ce{CH3COOH} et l'acide conjugué de \ce{NH3}.
  \item L'acide oxalique \ce{H2C2O4} et l'ion oxalate \ce{C2O4^{2-}} appartiennent-ils au même couple ?
        Justifier par le nombre de protons qui les sépare.
\end{enumerate}
\end{exobox}

\end{document}
